% s13_evidence.tex -- executed computations, counts, review record, failure semantics, missingness.
% Exported label: sec:evidence. Internal labels carry the prefix num:.
% Every number set by a macro \Vthree... or a framework macro comes from paper/generated/, which code/make_tables.py
% writes from the strict-JSON receipts under evidence/; code/check_consistency.py verifies the correspondence.
% The two lines below load the generated macros when the section is compiled on its own; in the full manuscript they are already loaded.
\ifdefined\VthreeExactPassed\else% Automatically generated by code/make_tables.py; do not edit.
% Values containing \times or a minus sign are for math mode.
\newcommand{\VthreeExactPassed}{148}
\newcommand{\VthreeExactAttempted}{148}
\newcommand{\VthreeFrameworkExact}{50}
\newcommand{\VthreeGaugeCycles}{66}
\newcommand{\VthreeGaugeCyclesPassed}{66}
\newcommand{\VthreeGaugeSaved}{31}
\newcommand{\VthreeGaugeExtended}{35}
\newcommand{\VthreeGaugeSamples}{137\,357}
\newcommand{\VthreeGaugeUnresolved}{0}
\newcommand{\VthreeGaugeExponentGap}{2.926\times10^{-11}}
\newcommand{\VthreeGaugeEtaZeroGap}{6.866\times10^{-11}}
\newcommand{\VthreeGaugePolarError}{8.341\times10^{-8}}
\newcommand{\VthreeGaugeMinSigma}{2.270\times10^{-3}}
\newcommand{\VthreeGaugeMinXi}{3.289\times10^{-6}}
\newcommand{\VthreeGaugeMinK}{1.133\times10^{-6}}
\newcommand{\VthreeGaugeMinW}{4.242\times10^{-7}}
\newcommand{\VthreeGaugeMaxLambda}{-1.021\times10^{-6}}
\newcommand{\VthreeGaugeMaxFormEigenvalue}{-1.308\times10^{-6}}
\newcommand{\VthreeGaugeUmin}{-24\,294}
\newcommand{\VthreeGaugeHolderLower}{0.999994}
\newcommand{\VthreeGaugeHolderUpper}{0.999990}
\newcommand{\VthreeGaugeEtaZeroNegative}{10}
\newcommand{\VthreeGaugeEtaOneNegative}{37}
\newcommand{\VthreeGaugeAttracting}{66}
\newcommand{\VthreeLgConstChecks}{141}
\newcommand{\VthreeLgConstAuthor}{55}
\newcommand{\VthreeLgConstReferee}{86}
\newcommand{\VthreeLgConstMinMargin}{1.2\times10^{-5}}
\newcommand{\VthreeLgLemmaChecks}{29\,561}
\newcommand{\VthreeLgCycleChecks}{474}
\newcommand{\VthreeLgBelowChecks}{96}
\newcommand{\VthreeLgSegmentChecks}{28\,991}
\newcommand{\VthreeLgMaxEtaRatio}{0.017}
\newcommand{\VthreeLgMinPeriodRatio}{1.304}
\newcommand{\VthreeLgMinStageRatio}{0.917}
\newcommand{\VthreeLgLogMtSixteen}{-13.98}
\newcommand{\VthreeLgLogMtSixty}{-67.5}
\newcommand{\VthreeLgLogMtThreeHundred}{-375.6}
\newcommand{\VthreeAttPairs}{13}
\newcommand{\VthreeAttBatteryPairs}{11}
\newcommand{\VthreeAttChecks}{156}
\newcommand{\VthreeAttOrderedPairs}{1.09\times10^{8}}
\newcommand{\VthreeAttNonnegPairs}{4}
\newcommand{\VthreeAttBeyondRounding}{0}
\newcommand{\VthreeAttWorstSeparated}{-2.0\times10^{-7}}
\newcommand{\VthreeAttGradPairs}{9}
\newcommand{\VthreeAttGradSamples}{2\,572}
\newcommand{\VthreeAttGradError}{5.7\times10^{-4}}
\newcommand{\VthreeAttGradMedian}{4.8\times10^{-7}}
\newcommand{\VthreeAttReturnError}{9.7\times10^{-10}}
\newcommand{\VthreeAttReturnFdError}{1.4\times10^{-4}}
\newcommand{\VthreeAttExponentGap}{1.8\times10^{-11}}
\newcommand{\VthreeAttCrestValueError}{6.8\times10^{-11}}
\newcommand{\VthreeWeakCycles}{20}
\newcommand{\VthreeWeakChecks}{180}
\newcommand{\VthreeWeakChordPairs}{5\,770\,615}
\newcommand{\VthreeWeakChordViolations}{0}
\newcommand{\VthreeWeakWindowError}{5.8\times10^{-6}}
\newcommand{\VthreeWeakTies}{4}
\newcommand{\VthreeWeakPhiYGood}{15}
\newcommand{\VthreeWeakExponentSpread}{1.1\times10^{-11}}
\newcommand{\VthreeWedgeChecks}{211}
\newcommand{\VthreeWedgeKurtosisToda}{1.7999988}
\newcommand{\VthreeWedgeKurtosisTodaMargin}{1.2\times10^{-6}}
\newcommand{\VthreeWedgeKurtosisOtherMargin}{1.5\times10^{-8}}
\newcommand{\VthreeWedgeMinRselPrime}{0.8000008}
\newcommand{\VthreeWedgeIdentityDefect}{3.2\times10^{-11}}
\newcommand{\VthreeWedgeIdentityCycles}{14}
\newcommand{\VthreeWedgeSelectionError}{4.0\times10^{-4}}
\newcommand{\VthreeWedgeExponentError}{7.5\times10^{-4}}
\newcommand{\VthreeWedgeRatioMin}{9.1}
\newcommand{\VthreeWedgeRatioMax}{9.9}
\newcommand{\VthreeWedgeRatioPredicted}{9.98}
\newcommand{\VthreeWedgeSelectionErrorAll}{1.8\times10^{-3}}
\newcommand{\VthreeWedgeExponentErrorAll}{6.7\times10^{-3}}
\newcommand{\VthreeCartChecks}{140}
\newcommand{\VthreeCartGridCycles}{300}
\newcommand{\VthreeCartBranchCycles}{827}
\newcommand{\VthreeCartHuntPoints}{240}
\newcommand{\VthreeCartHuntDirect}{239}
\newcommand{\VthreeCartMaxLambda}{-1.93\times10^{-6}}
\newcommand{\VthreeCartMaxLogMt}{-2.83\times10^{-4}}
\newcommand{\VthreeCartMinLogMt}{-505.2}
\newcommand{\VthreeCartBranchMaxLambda}{-2.0\times10^{-10}}
\newcommand{\VthreeCartHuntMaxLogMt}{-4.17\times10^{-4}}
\newcommand{\VthreeCartHuntPairs}{239\,103}
\newcommand{\VthreeCartHuntLongRuns}{1\,673}
\newcommand{\VthreeCartHuntNotFinished}{25}
\newcommand{\VthreeCartHuntNewton}{2\,868}
\newcommand{\VthreeCartRecomputedCycles}{51}
\newcommand{\VthreeCartRecomputedHunts}{24}
\newcommand{\VthreeCartRecomputedDiffT}{7.1\times10^{-12}}
\newcommand{\VthreeCartRecomputedDiffLogMt}{1.0\times10^{-9}}
\newcommand{\VthreeCartRecomputedHuntPairs}{24\,244}
\newcommand{\VthreeCartRecomputedLongRuns}{168}
\newcommand{\VthreeCartRecomputedNotFinished}{3}
\newcommand{\VthreeCartRecomputedNewton}{288}
\newcommand{\VthreeCartRecomputedNewPoints}{0}
\newcommand{\VthreeCartIndependentChecks}{23}
\newcommand{\VthreeCartIndependentDiffLogMt}{2.2\times10^{-9}}
\newcommand{\VthreeCartNegativeEhatCycles}{126}
\newcommand{\VthreeNegChecks}{123}
\newcommand{\VthreeNegRise}{0.0250}
\newcommand{\VthreeNegRatioAtE}{0.2871}
\newcommand{\VthreeNegSlope}{-0.132}
\newcommand{\VthreeNegSlopeMid}{0.052}
\newcommand{\VthreeNegSignChanges}{72}
\newcommand{\VthreeControls}{6}
\newcommand{\VthreeControlsPassed}{6}
\newcommand{\VthreeTotalAttempted}{30\,732}
\newcommand{\VthreeTotalPassed}{30\,732}
\newcommand{\VthreeRuntimeMinutes}{7.3}
\newcommand{\VthreePeakMemoryMb}{406}
\newcommand{\VthreeCartTopRhoMaxLogMt}{-251.0}
\newcommand{\VthreeCartTopRhoMinLogMt}{-505.2}
\fi
\ifdefined\ExponentAgreement\else% Automatically generated; do not edit.
\newcommand{\ExponentAgreement}{3.527\times10^{-12}}
\newcommand{\OrthogonalityError}{9.847\times10^{-11}}
\newcommand{\ReturnClosure}{2.098\times10^{-11}}
\newcommand{\MomentError}{1.159\times10^{-11}}
\newcommand{\CrossMatrixError}{2.443\times10^{-10}}
\newcommand{\CrossExponentError}{5.862\times10^{-14}}
\newcommand{\ReciprocalError}{1.290\times10^{-11}}
\newcommand{\EnvelopeHessianError}{1.807\times10^{-9}}
\newcommand{\EnvelopeDefectError}{8.046\times10^{-12}}
\fi
\section{Executed computations and what they establish}\label{sec:evidence}
The theorems of this paper rest on the written proofs. Computation enters a proof in one place only: the proof of Theorem~\ref{thm:largegain} uses a finite list of inequalities between explicitly computable numbers, and that list is recomputed here at 40 and at 50 significant digits. Every other computation in this section is a test. Exact symbolic computation tests the algebra of the identities. Floating-point computation on approximate orbits tests the proved statements for errors and supplies the evidence, for and against, concerning statements that are not proved. This section states what was executed, what each class of test can and cannot establish, the computed map of the weak exponent, how the proofs were reviewed, how failures are recorded, and what was not done.

\subsection{What was executed}\label{num:sec:executed}
The main command, \texttt{python code/run\_all.py}, re-runs the original suite; separate revision checks use Python script \path{code/rational_interval_audit.py} and \path{code/revision_controls.py}; Appendix~\ref{app:reproduce} gives the environment and the options. The run uses no network access and a single thread. On the machine described there it took $\VthreeRuntimeMinutes$ minutes, and the largest resident memory of any step was $\VthreePeakMemoryMb$\,MB. Each check writes a receipt in strict JSON under \path{evidence/}. A receipt records its name, its kind, whether it passed, the numbers of attempted and passed tests with the unit that is counted, the worst residuals, the parameters, the running time, and a statement of scope; every receipt that is not an exact symbolic equality carries the field \texttt{not\_interval\_certified}. Table~\ref{num:tab:counts} is generated from the receipts.

\begin{table}[ht]
\centering\footnotesize
\caption{Executed checks. The upper block accompanies Sections~\ref{sec:adjoint}--\ref{sec:geometry}, \ref{sec:compact}, \ref{sec:obstruction} and Appendix~\ref{app:hopf}; the lower block accompanies Sections~\ref{sec:gauges}--\ref{sec:attractor} and \ref{sec:largegain}--\ref{sec:negative}. The units differ from row to row, and the rows are not independent: several reuse the same computed cycles. No count is a probability that a theorem is correct.}\label{num:tab:counts}
% Automatically generated by code/make_tables.py; do not edit.
\begin{tabular}{@{}>{\raggedright\arraybackslash}p{5.3cm}>{\raggedright\arraybackslash}p{2.7cm}rr>{\raggedright\arraybackslash}p{3.5cm}@{}}
\toprule
Check & Kind & Passed & Attempted & Unit counted \\
\midrule
Framework identities & exact & 50 & 50 & residual tested equal to zero \\
Periodic boundary-value solves & floating point & 31 & 31 & cycle \\
Cycle geometry and adjoint audits & floating point & 31 & 31 & cycle \\
DOP853--Radau repetitions & floating point & 4 & 4 & cycle \\
Reciprocal-rate comparisons & floating point & 12 & 12 & pair of cycles \\
Envelope Hessian checks & floating point & 4 & 4 & phase \\
Deliberate-error controls (framework) & control & 4 & 4 & wrong assertion rejected \\
\midrule
New algebraic identities & exact & 148 & 148 & residual tested equal to zero \\
Gauges and exponent formulas on cycles & floating point & 66 & 66 & cycle \\
Constants of the large-gain proof & high precision & 141 & 141 & scalar comparison \\
Large-gain lemmas on cycles and on non-periodic segments & floating point & 29\,561 & 29\,561 & inequality evaluated \\
Attractor: supporting planes, gradient, crest-axis formula & floating point & 156 & 156 & parameter pair and predicate \\
Weak adjoint: event order, dual planes, rank-one identity & floating point & 180 & 180 & cycle and predicate \\
Kurtosis bound, period identities, third-order formula & high precision and floating point & 211 & 211 & level, identity, cycle or table row \\
Cartography: shipped tables and recomputed subsample & floating point & 140 & 140 & table statement, cycle or search \\
Refutations reproduced & floating point & 123 & 123 & refutation, identity or candidate \\
Deliberate-error controls & control & 6 & 6 & wrong assertion rejected \\
\bottomrule
\end{tabular}

\end{table}

The upper block of the table is unchanged in content from the earlier release of this work. It consists of $\VthreeFrameworkExact$ exact identities; 31 cycles computed as periodic boundary-value problems with prescribed crest amplitude $h=U_{\max}$, for $\mu\in\{0.1,0.3,1,10/3,10\}$ and $h\in\{0.05,0.2,0.5,1,2,3\}$ together with $(\mu,h)=(0.3,0.01)$; an audit of each of these cycles; repetitions with an implicit integrator; comparisons between reciprocal rates; and checks of the local Hessian of the supporting envelope. The boundary-value problem is posed in the time $\tau=\sqrt\mu\,t$ for the state $(U,p/\sqrt\mu,y/k)$, with periodicity, $U(0)=h$ and $p(0)=0$ imposed and the period and the gain as unknowns; the collocation tolerance is $2\times10^{-9}$ ($2\times10^{-11}$ for the smallest cycle), exponentials are not clipped, and convergence of the solver and scientific acceptance are separate fields of the receipt. The default run reuses the 31 saved solutions; the option \texttt{-{}-regenerate} recomputes them. On each saved cycle the audit integrates the adjoint ratios, and separately a normalised positive solution of the linear adjoint system, and separately again the nonlinear flow with two variations of the crest section. The largest difference between the two values of the weak exponent is $\ExponentAgreement$, the largest scaled defect of tangent orthogonality is $\OrthogonalityError$, the largest closure error of the return is $\ReturnClosure$ and the largest relative residual of the moment identities is $\MomentError$. All 31 cycles show the six extrema in the order of Theorem~\ref{thm:six} and winding number $-1$. Four repetitions with the Radau method change the return matrix by at most $\CrossMatrixError$ and the exponent by at most $\CrossExponentError$; twelve comparisons between the rates $\mu$ and $1/\mu$ agree to $\ReciprocalError$; the envelope Hessian agrees with its finite-difference value to $\EnvelopeHessianError$ in the scaled norm.

\subsection{Three kinds of check}\label{num:sec:classes}
\paragraph{Exact symbolic equalities.}
A test of this kind forms a polynomial, or the numerator of a rational expression after explicit denominators are cleared, expands it with the computer algebra system, and compares it with the zero polynomial; only tests of the form ``residual equals zero'' are counted. The flow is used with $x=ke^U$ as an independent symbol obeying $x'=xp$, which makes every identity polynomial. An identity that holds only on the constraint $O=\zeta p-\eta p'-(x-y)=0$ of tangent orthogonality is tested in the form ``residual equals cofactor times $O$'', with the cofactor obtained by exact division and written into the receipt. The $\VthreeExactAttempted$ identities of the lower block cover: the gauge equations \eqref{eq:gaugeodes} with their cofactors ($-U$ for $\Sg$, and $\eta p+k-\mu p+p-y$ for $\Xg$), the state forms of $\Sg'$ and $\Xg'$, the relation $\eta-\etaz=-pK/\Sg$, and the shifted-centre gauge; the scalar Riccati equation of Proposition~\ref{prop:kappa} in its three forms; the equation of $\LamE$ and the decomposition of $\Sg$ at the equilibrium; the polar form of Proposition~\ref{prop:polar} and the difference formula of Proposition~\ref{prop:nestedradii}; the Lotka--Volterra chain and the second-order equations of Proposition~\ref{prop:lvchain}; the equations of the strong-stable ratios, the weak adjoint and its components, the three crossing derivatives used in the proof of Theorem~\ref{thm:twelve}, the chord and curvature pairings behind Theorem~\ref{thm:dualsupport}, and the carrier identity of Proposition~\ref{prop:rankone}; the rows of the compound system for the wedges that occur in Sections~\ref{sec:gauges} and \ref{sec:negative}; the first-order coefficients of Lemma~\ref{lem:rigidity}; the period identities of Lemma~\ref{lem:tauidentities}, each as the exact derivative of an explicit primitive, and the two polynomial identities in the proof of Theorem~\ref{thm:kurtosis}; the Lyapunov identity and the cooperative system of Lemma~\ref{lem:lglinear}; and the wedge, divergence and two-orbit identities of Section~\ref{sec:attractor}. All $\VthreeExactPassed$ pass. Such a test establishes a finite algebraic identity and nothing else. It does not establish positivity of a gauge, an averaging step, a limit, a Perron--Frobenius or Sturm argument, or any topological step; these are in the written proofs.

\paragraph{High-precision constants.}
The constants of Section~\ref{sec:largegain} are evaluated by two scripts that were written independently of one another, one with the proof and one from the text of the proof alone, at 40 and at 50 significant digits. Together they make $\VthreeLgConstChecks$ comparisons, $\VthreeLgConstAuthor$ in the first script and $\VthreeLgConstReferee$ in the second. The second script also compares the closed forms of $c_{\rm i}$, $c_{\rm ii}$ and $\Psi_\mu$ with quadrature, tests on grids the monotonicity statements by which every side condition is reduced to $\lambda=k_0$ (601 values of $\lambda$ between $10^3$ and $10^{15}$), and tests the exclusion step of the bootstrap on a grid of $6\times5\times2001$ points. All comparisons pass. The smallest relative margin among the scalar comparisons is $\VthreeLgConstMinMargin$, against a working precision of $10^{-40}$; one grid test contains cases of equality at $\mu=1$ and is made with an allowance of $10^{-40}$. Those two scripts use non-interval arithmetic and finite grids. A separate v3.1 rational-interval audit encloses selected endpoints, including $3/4+q(1000)<0.83$, but not the grids or analytic reductions. Table~\ref{num:tab:lgconstants} lists the recomputed values in the notation of Table~\ref{lg:tab:constants}.

\begin{table}[ht]
\centering\small
\caption{Constants of the proof of Theorem~\ref{thm:largegain}, recomputed at 40 significant digits and rounded (receipt \texttt{evidence/v3/v3\_largegain\_constants.json}).}\label{num:tab:lgconstants}
% Automatically generated by code/make_tables.py from evidence/v3/v3_largegain_constants.json; do not edit.
\begin{tabular}{@{}lll@{}}
\toprule
Quantity & Recomputed value & Compared with \\
\midrule
$\omega_0(k_0)$, $T_0(k_0)$ & $8.0467$, $5.89277$ & $T_0(k_0)\ge a_1+a_2+a_3+a_4=5.89$ \\
$k_0c_{\rm i}(a_1)$, $\chi_{\rm ii}(a_2)$, $k_0c_{\rm b}(a_4)$ & $1388.7$, $102.4$, $2.369$ & $\ge2$, $\ge2$, $\ge L_\ast/\varphi_\ast=2.313$ \\
bound for $c_1$ & $0.37453$ & $\ge c_1=0.374$ \\
$\tilde R(k_0)$, $\bar\varepsilon_1$, $\vartheta_1(k_0)$, $T_1(k_0)$ & $14.921$, $0.2151$, $0.19402$, $43.898$ & $\vartheta_1<1$, $T_1\ge22$ \\
$z_{10}$, $\chi_0$ & $10.169$, $0.20969$ &  \\
bound for $c_p$; $c_T$ & $0.81943$; $0.18333$ & $\ge c_p=0.819$ \\
$\hat R$, $\hat L$, $\hat\varepsilon_1$ at $\lambda=k_0$ & $13.8175$, $7.8886$, $0.06637$ &  \\
$\hat\vartheta_N$, $a_c$, $A_1$ at $\lambda=k_0$ & $0.07772$, $3.6345$, $23.1755$ & $\hat\vartheta_N<1$ \\
$\bar\tau_1$, $\hat w$, $v_f$ at $\lambda=k_0$ & $0.028297$, $28135$, $0.040066$ &  \\
$C_\tau(k_0)$, $c_{\log}$ & $71.785$, $12.272$ & $\le71.79$, $\le12.28$ \\
$q(k_0)$ & $0.077177$ & $<0.0772$ \\
final margin & $45.83-5.54-8.61=31.68$ & $>0$ \\
\bottomrule
\end{tabular}

\end{table}

\paragraph{Floating-point diagnostics.}
The remaining checks of the lower block integrate the differential equations in double precision with an explicit Runge--Kutta method of order eight, at relative tolerances between $10^{-10}$ and $10^{-13}$. A cycle is located by Newton iteration on the crest return map, with the derivative from the variational equations or from finite differences, and at large gain by forward iteration. The troughs of the cycles are deep: on the cycles of the gauge check $U$ falls to $\VthreeGaugeUmin$, and on the large-gain cycles and in the cartography much lower, so that $x=ke^U$ underflows. The state and the adjoint ratios are therefore integrated in the logarithmic variables $U$, $\log(\lambda-p)$, $\log y$, $\log\eta$, $\log\zeta$, which keep relative accuracy. Every sign test is evaluated as a sum of signed monomials with stored logarithms of their magnitudes and is reported as the ratio of the sum to the sum of the absolute values of its terms; a sign counts as resolved if this ratio exceeds $10^{-9}$ in modulus. For $\mu>1$ the gauge $W$ cancels to leading order in the trough in these variables; it is then evaluated from an identity in the variables $\eta-(\mu-1)$ and $y-(\mu-1)(\lambda-p)$, which are positive and are integrated in logarithmic form. Predicates and tolerances are fixed in the scripts and written into the receipts. A check of this kind can detect an error in a formula or in a sign, and it supplies numerical values. It cannot establish a statement about every cycle, because at each parameter pair the computation produces one cycle; positivity at sampled points is not positivity on a continuum; and no error bound is enclosed.

\subsection{Results for the proved statements}\label{num:sec:results}
\paragraph{Gauges and exponent formulas (Section~\ref{sec:gauges}).}
The cycles are the 31 saved cycles and $\VthreeGaugeExtended$ further cycles at $\mu\in\{0.02,0.1,0.5,1,3,10,50\}$, $\rho\in\{1.001,1.1,2,6,20\}$, in total $\VthreeGaugeSamples$ sample points. On all $\VthreeGaugeCycles$ cycles the gauges $\Sg$, $\Xg$, $K$ and $W$ are positive at every sample, with no unresolved sign; the smallest ratios are $\VthreeGaugeMinSigma$ for $\Sg$, $\VthreeGaugeMinXi$ for $\Xg$, $\VthreeGaugeMinK$ for $K$ and $\VthreeGaugeMinW$ for $W$. The three expressions $\avg\eta-\mu$, $-\avg{pK/\Sg}$ and $-\avg{pW/\Xg}$ of \eqref{eq:exactexponent} agree to $\VthreeGaugeExponentGap$ relative to the mean absolute value of their integrands; $\avg{\etaz}=\mu$ holds to $\VthreeGaugeEtaZeroGap$ in the same sense; the polar angle of Proposition~\ref{prop:polar} advances by $2\pi$ per period to $\VthreeGaugePolarError$; the quadratic form of Proposition~\ref{prop:monotoneintegrals} is negative definite, with largest eigenvalue $\VthreeGaugeMaxFormEigenvalue$; and the two inequalities of Proposition~\ref{prop:holder} hold, the ratios of the smaller to the larger side being at most $\VthreeGaugeHolderLower$ and $\VthreeGaugeHolderUpper$ (the cycles closest to the Hopf curve). All $\VthreeGaugeAttracting$ cycles have $\lambda_c<0$; the largest value is $\VthreeGaugeMaxLambda$. The state functions $\etaz$ and $\eta_1$ of Section~\ref{sec:gauges} are not positive in general: $\eta_1$ is negative on part of $\VthreeGaugeEtaOneNegative$ of the $\VthreeGaugeCycles$ cycles, and $\etaz$ on part of $\VthreeGaugeEtaZeroNegative$ of them, all with $\mu\le0.1$; on the cycles of small amplitude both are positive throughout.

\paragraph{Weak adjoint (Section~\ref{sec:weakadjoint}).}
On $\VthreeWeakCycles$ cycles, $\mu\in\{0.02,0.3,1\}$ with $\rho\in\{1.01,1.5,4,12\}$ and eight further pairs with $\mu$ up to $50$, each of $p$, $p'$, $y'$ and the three components of the weak adjoint changes sign exactly twice per period, in the cyclic order of Theorem~\ref{thm:twelve}. On $\VthreeWeakTies$ cycles with deep troughs the zero of $p$ at the trough and the adjacent zero of the second component fall in one double-precision time cell; their order is then decided from the exact relation between that component and $y'$ at a zero of $p$. The determinants of Theorem~\ref{thm:dualsupport} are negative at all $\VthreeWeakChordPairs$ pairs of base point and chord examined, at the equilibrium, and for the curvature vector. The identity of Proposition~\ref{prop:rankone} is reproduced with relative error at most $\VthreeWeakWindowError$ for the test functions $1$, $p'$, $-U$, $\cos\omega t$ and $\sin3\omega t$, with $\lambda_c$ taken from the independent forward integration of $(\eta,\zeta)$. For the test function $y$ the comparison is reported and not counted: on five cycles with deep troughs and $\mu<1$ the right-hand side is many orders of magnitude smaller than the two integrals whose difference it is, and the comparison is rounding noise; on the other $\VthreeWeakPhiYGood$ cycles it agrees to $10^{-5}$. The three values of $\lambda_c$, from $\avg\eta-\mu$, from $\avg{\mathfrak r}-A$ and from the window identity, agree to $\VthreeWeakExponentSpread$. The inequality $\mathfrak r>\max(1,\mu)$ of Lemma~\ref{lem:strongstable} cannot be tested by this computation, because the margin is integrated in logarithmic form and is positive by construction; the computation confirms only that the two formulations of the margin agree where both are resolved.

\paragraph{Attractor (Section~\ref{sec:attractor}).}
At $\VthreeAttBatteryPairs$ parameter pairs with $\mu\in\{0.02,0.3,1,3\}$ and $\rho$ between $1.01$ and $12$ the attractor is sampled by orbits of the unstable manifold of $E$ at several phases together with the cycle. The supporting-plane inequality of Theorem~\ref{thm:twoorbit} is evaluated for every ordered pair of samples; at two further pairs, $(\mu,\rho)=(3,12)$ and $(10,4)$, where the trough is deep and $\mu>1$, it is evaluated in the variables without cancellation described above. In total these are $\VthreeAttOrderedPairs$ ordered pairs. No pair violates the inequality beyond the rounding amplification of the difference of two nearly coincident points, and among pairs with relative separation above $10^{-6}$ the largest value of the normalised left-hand side is $\VthreeAttWorstSeparated$; $\VthreeAttNonnegPairs$ pairs, all closer than that, give a non-negative value below their rounding threshold. $K$ and $W$ are positive on all samples of the spirals. The relation $\nabla h=(\zeta,-\eta)$ is tested by finite differences across neighbouring orbits of the unstable manifold: at the $\VthreeAttGradPairs$ pairs where neighbouring orbits are transversally separated the relative error is at most $\VthreeAttGradError$ on $\VthreeAttGradSamples$ samples, with median at most $\VthreeAttGradMedian$; the truncation error of the stencil is about $2\times10^{-3}$. Close to the Hopf curve neighbouring orbits differ almost by a time shift and the test is not applicable; this is so at $(0.3,1.01)$ and $(1,1.1)$. The formula \eqref{att:eq:returnderivative} for the derivative of the return map on the crest half-axis agrees with the derivative obtained from the variational equations of the three-dimensional flow to $\VthreeAttReturnError$, and with finite differences between neighbouring orbits to $\VthreeAttReturnFdError$. The exponent $\avg\eta-\mu$ agrees to $\VthreeAttExponentGap$ with the sum of the two leading exponents from a discrete QR computation that does not use the Riccati system. These tests concern the formulas of Theorem~\ref{thm:attractor} and Proposition~\ref{prop:crestaxis}; they do not test the topological steps of the proof.

\paragraph{Large gain (Section~\ref{sec:largegain}).}
The inequalities of the lemmas are evaluated on computed cycles with $k\in\{1000,3000\}$ and $\mu$ between $0.02$ and $1$, by two independent integrators ($\VthreeLgCycleChecks$ evaluations), and their derivation steps on orbit segments that are not periodic: random negative lobes, random spikes, random lobes satisfying the hypotheses of the two stages, and, for Lemma~\ref{lem:lglinear}, linear cooperative systems with an arbitrary periodic coefficient ($\VthreeLgSegmentChecks$ evaluations). All pass. The computed cycles are far inside the proved bounds: $T\mu/\lambda\ge\VthreeLgMinPeriodRatio$ against $c_T=0.1833$, and $\avg\eta/\mu\le\VthreeLgMaxEtaRatio$ against $0.83$; on the non-periodic lobes the smallest value of $p_k/k$ under the hypotheses of the first stage is $\VthreeLgMinStageRatio$ against $c_1=0.374$. Below the threshold, at $\mu=1$ and $k=16$, $60$, $300$, the cycles have $\log m_t=\VthreeLgLogMtSixteen$, $\VthreeLgLogMtSixty$ and $\VthreeLgLogMtThreeHundred$. A cycle is a single orbit, so the cycle tests do not probe the quantifier over all cycles; the tests on non-periodic segments exercise the steps of the derivation in regimes that the computed cycle does not visit.

\paragraph{Small rate ratio (Section~\ref{sec:wedge}).}
The inequality of Theorem~\ref{thm:kurtosis} is evaluated by quadrature at 30 significant digits on 14 levels $\xi\in[10^{-6},10^{6}]$ of the potential $e^U-U-1$ and on ten further convex potentials: the ratio $\avg{P^4}/\avg{P^2}^2$ stays below $9/5$, by at least $\VthreeWedgeKurtosisTodaMargin$ and $\VthreeWedgeKurtosisOtherMargin$ respectively, and the smallest value of $R_{\mathrm{sel}}'$ on the grid is $\VthreeWedgeMinRselPrime$. Two non-convex potentials violate the bound, which shows that the test can fail. The identities of Lemma~\ref{lem:tauidentities} hold on $\VthreeWedgeIdentityCycles$ computed cycles with normalised defect at most $\VthreeWedgeIdentityDefect$. For cycles of prescribed crest amplitude $0.3$, $1$, $2$, $3$ and $\mu$ from $3\times10^{-2}$ down to $10^{-4}$ the third-order predictions of Theorem~\ref{thm:wedge} are approached with errors that decrease by a factor between $\VthreeWedgeRatioMin$ and $\VthreeWedgeRatioMax$ from $\mu=10^{-3}$ to $\mu=10^{-4}$, where second order in $\varepsilon$ predicts $\VthreeWedgeRatioPredicted$; at $\mu=10^{-4}$ and amplitude at most $1$ the relative errors are $\VthreeWedgeSelectionError$ for the selection rule and $\VthreeWedgeExponentError$ for $\log m_t$, and at amplitude $3$ they are $\VthreeWedgeSelectionErrorAll$ and $\VthreeWedgeExponentErrorAll$. The rows of Tables~\ref{wed:tab:selection} and \ref{wed:tab:fixedrho} are recomputed and agree with the printed values to the printed digits. The agreement is evidence for the asymptotic formulas; Theorem~\ref{thm:wedge} uses the source-mapped slow-graph lemma, not these computations.

\paragraph{Refutations (Section~\ref{sec:negative}).}
A candidate statement counts as refuted by floating-point computation if a violation is found that is far above rounding and stable under a change of tolerance. Four of the refutations of Section~\ref{sec:negative} are reproduced in the run at their stated parameters. At $(\mu,\rho)=(0.0027,5.02)$ the logarithm of $R_1(u)/u$ on the crest half-axis equals $\VthreeNegRatioAtE$ at the equilibrium, has an interior minimum, and rises again by $\VthreeNegRise$; hence $R_1(u)/u$ is not monotone and $R_1$ is not concave there. At $(1,5)$ the ratio $(\mu-\zeta(u,0))/(k(e^u-1))$ has slope $\VthreeNegSlope$ just inside the cycle and $\VthreeNegSlopeMid$ at middle amplitudes, so that it is not monotone. On 18 cycles, $\mu\in\{0.02,0.1,0.5,1,3,10\}$ and $\rho\in\{1.1,2,6\}$, each of the four pointwise candidates $x/\zeta-A$, $\zeta/\eta-1$, $k\psiU(U)/K-A$ and $xp^2/W-A$ takes both signs, in all $\VthreeNegSignChanges$ cases, although its mean is $\lambda_c<0$; and $K-(\zeta-\mu)\varkappa(U)$ takes both signs at $\rho=6$, $\mu\le1$. The three averages used in the proof of Proposition~\ref{prop:nomonomial} are confirmed on the same cycles, and the identities of Proposition~\ref{wed:prop:melnikov} on five levels of the centre. The other entries of Table~\ref{neg:tab:candidates} rest on outputs of the research phase, which are shipped but not re-executed; see Section~\ref{num:sec:missing}.

\subsection{Cartography of the weak exponent}\label{num:sec:cartography}
The tables under \path{evidence/v3/cartography/} were computed in the research phase with a compiled integrator and are shipped unchanged. A grid of $\VthreeCartGridCycles$ cycles covers the twelve values $0.002$, $0.005$, $0.01$, $0.02$, $0.05$, $0.1$, $0.2$, $0.3$, $0.5$, $0.7$, $0.85$, $1$ of $\mu$ and $25$ values of $\rho-1$, logarithmically spaced from $10^{-3}$ to $199$; each cycle was reached by continuation in the crest amplitude from the Hopf point, which gives $\VthreeCartBranchCycles$ further cycles on the twelve branches. By the symmetry $\mu\mapsto1/\mu$ at fixed $\rho$ recalled in Section~\ref{sec:model}, which transports cycles and their multipliers, the grid also describes the reciprocal rates. Table~\ref{num:tab:cartography} shows part of it. The tables support the following statements, each of which is checked by the run.
\begin{enumerate}[label=\textup{(\roman*)}]
\item Every computed cycle has $\lambda_c<0$. The largest value on the grid is $\VthreeCartMaxLambda$, at $\mu=0.002$, $\rho-1=10^{-3}$, where $\log m_t=\VthreeCartMaxLogMt$; on the branches the largest value is $\VthreeCartBranchMaxLambda$, at the smallest amplitude.
\item At each $\mu$, $\log m_t$ and $\lambda_c$ are strictly decreasing in $\rho$ on the grid; at $\rho=200$, $\log m_t$ lies between $\VthreeCartTopRhoMinLogMt$ and $\VthreeCartTopRhoMaxLogMt$.
\item On each of the twelve branches $\rho$ is strictly increasing in the crest amplitude. A fold of the branch would appear as a failure of this monotonicity.
\item No second cycle was found. At $\VthreeCartHuntPoints$ random parameter pairs, with $\mu$ log-uniform between $10^{-3}$ and $10^{3}$ and $\rho-1$ log-uniform between $10^{-3}$ and $299$, three searches were made: $\VthreeCartHuntPairs$ pairs $(U_n,U_{n+1})$ of successive crests on the attractor, inside and outside the cycle, none with the displacement $U_{n+1}-U_n$ of the sign opposite to that of $U_\ast-U_n$, where $U_\ast$ is the crest of the located cycle; $\VthreeCartHuntLongRuns$ long runs from the unstable manifold of $E$, from crests far outside, and from random initial states, none of which converged to a different crest ($\VthreeCartHuntNotFinished$ did not finish, being too slow near the Hopf curve or too costly); and $\VthreeCartHuntNewton$ Newton iterations on the crest map from perturbed guesses, which found no new fixed point. At one of the $\VthreeCartHuntPoints$ points the integrator failed in the direct coordinates, and the point was computed at the reciprocal rate.
\end{enumerate}

\begin{table}[ht]
\centering\small\setlength{\tabcolsep}{4.5pt}
\caption{$\log m_t=\lambda_cT$ on part of the grid (floating point; from \texttt{grid\_table.csv}).}\label{num:tab:cartography}
% Automatically generated by code/make_tables.py from evidence/v3/cartography/grid_table.csv; do not edit.
\begin{tabular}{@{}lrrrrrr@{}}
\toprule
$\rho-1$ & $\mu=0.002$ & $\mu=0.01$ & $\mu=0.05$ & $\mu=0.2$ & $\mu=0.5$ & $\mu=1$ \\
\midrule
$0.001$ & $-2.83\times10^{-4}$ & $-6.16\times10^{-4}$ & $-1.28\times10^{-3}$ & $-2.05\times10^{-3}$ & $-2.42\times10^{-3}$ & $-2.51\times10^{-3}$ \\
$0.0127$ & $-4.17\times10^{-3}$ & $-7.83\times10^{-3}$ & $-1.60\times10^{-2}$ & $-2.59\times10^{-2}$ & $-3.07\times10^{-2}$ & $-3.19\times10^{-2}$ \\
$0.161$ & $-0.103$ & $-0.116$ & $-0.171$ & $-0.292$ & $-0.373$ & $-0.396$ \\
$1.23$ & $-1.02$ & $-1.04$ & $-1.14$ & $-1.47$ & $-2.00$ & $-2.31$ \\
$9.42$ & $-9.20$ & $-9.30$ & $-9.79$ & $-11.6$ & $-15.2$ & $-19.5$ \\
$199$ & $-251$ & $-253$ & $-263$ & $-302$ & $-380$ & $-505$ \\
\bottomrule
\end{tabular}

\end{table}

The run recomputes a subsample fixed in advance. With the integrator of this bundle, in logarithmic variables and started from the shipped crest, $\VthreeCartRecomputedCycles$ grid cycles are recomputed: four values of $\rho$ at every $\mu$, and $\rho=200$ at $\mu\in\{0.2,0.5,1\}$. The periods agree with the table to $\VthreeCartRecomputedDiffT$ and $\log m_t$ to $\VthreeCartRecomputedDiffLogMt$, relative to $1+|\log m_t|$. The searches are repeated at $\VthreeCartRecomputedHunts$ of the $\VthreeCartHuntPoints$ points with the compiled code: $\VthreeCartRecomputedHuntPairs$ pairs with no wrong sign, $\VthreeCartRecomputedLongRuns$ long runs of which $\VthreeCartRecomputedNotFinished$ did not finish and none ended elsewhere, and $\VthreeCartRecomputedNewton$ Newton iterations with $\VthreeCartRecomputedNewPoints$ new fixed points. The located cycle of $\VthreeCartIndependentChecks$ of these points is confirmed by the integrator of this bundle to $\VthreeCartIndependentDiffLogMt$ in $\log m_t$. The searches are limited by a time budget, so their counts vary from run to run.

The cartography is evidence and proves nothing. It is finite; it covers $\mu$ between $10^{-3}$ and $10^{3}$ and $\rho\le300$; outside the cycle the return map on the attractor is defined only for orbits whose transient has decayed; and a family of cycles lying entirely outside the scanned range of amplitudes, or between two scanned seeds, would not be seen.

\subsection{Record of adversarial review}\label{num:sec:review}
The proofs have not been refereed by specialists. The review that took place was internal and adversarial, and it is recorded in full in the bundle. Each new result was first written, with its proof, in a working document of one research line. An independent referee was then instructed to refute the statements marked as proved in that document: to re-derive every step, to look for circular use of stability or uniqueness, and to search numerically for counterexamples with scripts that share no code with those of the author. The authors of the working documents and the referees were instances of a large language model working independently of one another; a model of a different developer was used for a further adversarial reading of the hardest open step and of Sections~\ref{sec:gauges}--\ref{sec:negative} as drafted for this paper. The reports are in \path{provenance/research_lanes/} (the file \texttt{VERIFY.md} of each line, and the directory \texttt{sol} for the second model), and the report file that accompanies each section lists the corrections applied to it. The outcome, statement by statement, was as follows.
\begin{itemize}
\item \emph{Gauge theorem and its corollaries} (Theorem~\ref{thm:gauges}, Propositions~\ref{prop:kappa}--\ref{prop:monotoneintegrals}). Obtained independently in three research lines and refereed in two. No gap was found in the proofs; one referee integrated 21 cycles at 400 bits and another 65 cycles with independent code. One by-product was faulted: the statement that a Dulac function cannot be regular at $E$ had been derived with an averaging step that needs a hyperbolic cycle, and with a growth clause argued only for power laws. Proposition~\ref{prop:nodulac} is the repaired statement.
\item \emph{Weak adjoint} (Lemma~\ref{lem:strongstable}, Theorems~\ref{thm:twelve} and \ref{thm:dualsupport}, Proposition~\ref{prop:rankone}). No gap. The working document called $\avg{\mathfrak r}<A$ equivalent to the conjecture; it is equivalent to $\lambda_c<0$, which implies uniqueness and is not known to be implied by it, and the paper says so. For the monotonicity of the crest map (Proposition~\ref{prop:crestmonotone}) the referee found that the existence and transversality of the first return on the whole section had not been argued, and supplied the argument; the second model then found two steps of the written proof too compressed, and they were written out.
\item \emph{Attractor theorem} (Theorems~\ref{thm:attractor} and \ref{thm:twoorbit}, Proposition~\ref{prop:crestaxis}). No gap after a line-by-line audit and tests at 15 parameter pairs. One identity outside the theorem was faulted: a sign statement for the radial scaling had been made without the restriction $0\le\theta\le1$ and is false for $\theta>1$; Proposition~\ref{att:prop:scaling} states the corrected form. The second model raised fourteen points on the drafted section, among them an imported lemma on transversals stated without the definition that makes it true and strictness claims in Proposition~\ref{prop:crestaxis}; their dispositions are recorded with the section, and the repaired text has not been read again by a referee.
\item \emph{Large-gain theorem} (Theorem~\ref{thm:largegain}, Corollaries~\ref{cor:lgunique} and \ref{cor:lgreciprocal}). A first reading by the second model found no error in the lemmas and one false sentence in the proof of the reciprocal corollary, which asserted a correspondence of positive domains; the proof was rewritten to transport only cycles and multipliers. The referee then re-derived every lemma and every constant, read the cited convergence theorem in its source, and found no gap; about $20\,000$ inequality evaluations on 88 cycles and $2300$ non-periodic segments gave no counterexample. A second reading of the section as ported found eight minor points, all applied.
\item \emph{Small rate ratio} (Lemma~\ref{lem:tauidentities}, Theorem~\ref{thm:kurtosis}, Theorem~\ref{thm:wedge}). The premise from which this line started, that the first-order Melnikov function has a zero, was refuted by the line itself (Proposition~\ref{wed:prop:melnikov}). The referee found the kurtosis theorem and the identities correct, and in the wedge theorem one regularity slip ($C^1$ written where $C^2$ is needed and available), which was repaired. The second model found the identities, Proposition~\ref{wed:prop:melnikov}, the kurtosis theorem and the normal-form lemma correct, and asked for details of the conditional theorem to be written out and for the status of the imported invariant-manifold theorem to be stated. The details were written out, and version 3.0 exposed the property as assumption (F). Version 3.1 replaces the unchecked sketch with Lemma~\ref{wed:lem:slowgraph}, a cutoff application of Eldering Theorem~3.2 checked in a targeted internal source-and-proof pass.
\item \emph{Compactness and the neutral locus} (Sections~\ref{sec:compact} and \ref{sec:obstruction}). The second model found the statements on the state box, the separation from the equilibrium, the speed and the compactness correct, and a gap in the Hopf-collar lemma (the orientation of a crossing, the wording of the centre-manifold step, and unchecked hypotheses of an imported step), which was repaired.
\item \emph{Negative results} (Section~\ref{sec:negative}). Propositions~\ref{prop:holder}, \ref{prop:lvchain}, \ref{prop:gaussmap} and \ref{prop:nomonomial} come from a research line that had no referee of its own; they were examined only by the second model, in the drafted sections. It found six statements of Section~\ref{sec:negative} correct, asked for two steps in the proof of Lemma~\ref{lem:rigidity} to be written out, and marked several sentences as claiming more than was shown; its findings are in the review directory.
\item \emph{The open sign.} The second model was asked to refute the assessment that $\lambda_c<0$ on every cycle cannot be obtained from the structure established here by an elementary argument. It returned a partial result: it proved only statements already in this paper and identified the same missing monotonicity of a one-turn integral as Proposition~\ref{neg:prop:signchange}.
\end{itemize}
Alarms that turned out to be defects of a checker are also recorded in the reports: apparent violations of the supporting-plane inequality and of $W>0$ in deep troughs when the adjoint ratios are integrated in the original variables, sign anomalies for pairs of points closer than the rounding resolution, a cancellation for $\mu>1$, and an equality case compared without tolerance. Each was traced, and the checks of this section use the formulations that avoid them. This record describes automated adversarial reading. It lowers the chance of an undetected error and is not a substitute for refereeing by specialists. The framework of Sections~\ref{sec:adjoint}--\ref{sec:geometry}, on which the new results depend, was assessed in one report on the earlier release, reproduced in \path{review/v2_0_referee_report.md}; the referees of the new results re-derived the framework statements that they import.

\subsection{Failure semantics}\label{num:sec:failure}
The run exits with a non-zero status if any required step fails or if any required receipt does not pass; every step is executed even after an earlier failure. A receipt passes only if every one of its predicates passes, a computation that cannot be carried out is recorded as a failure, and a test declared not applicable is listed under that name and is not counted. A final step verifies that every number printed in this section equals the value in its receipt, that every receipt and file named in the paper exists, and that the receipts were produced by the shipped code.

Deliberately wrong assertions must fail. The four controls of the upper block and the $\VthreeControls$ controls of the lower block pass only when the wrong assertion is rejected by the same predicate and exit status as in the real check: the star-gauge identity with the wrong sign of $pK$ (rejected as one non-zero residual among 149); the assertion $\Sg>0$ on a non-periodic orbit started at $(U,p,y)=(0.5,0,k/2)$, where it is false; the large-gain constants $c_1=0.38$ and $c_p=0.83$ in place of $0.374$ and $0.819$; the exponent formula with the sign of $\avg{pK/\Sg}$ reversed; and the twelve events in reversed order. The receipts of the rejected assertions are kept and are marked as failing.

Three genuine failures of checkers are kept with their receipts. First, in the earlier release one exact check failed (49 of 50) because the off-constraint form of an identity used the wrong multiple of $O$; the identity on the constraint was unaffected. Second, the first version of the gauge check failed on one of 66 cycles, at $(\mu,\rho)=(0.02,20)$: the residual of the located crest was scaled by $k$ although the crest value of $y$ is $775k$ there, and the scaled residual $1.8\times10^{-9}$ exceeded the gate $10^{-9}$ while every mathematical predicate passed; the residual is now measured relative to the crest state. In the same version the sign of $W$ was unresolved at $11\,236$ of $137\,303$ samples, all at $\mu\ge3$, which led to the cancellation-free evaluation described above. Third, the lemma script written with the large-gain proof, run at $(\mu,k)=(0.02,1000)$, passes 36 of its 37 comparisons. The remaining one, $M_+<kT_-$, compares two numbers near $3.26\times10^{9}$ that differ by $3\times10^{-13}$ in relative terms, while the script obtains $M_+$ by a trapezoid rule whose own relative error is $7\times10^{-12}$; the comparison is unresolved and is recorded as failed. That run takes seven minutes and is not part of the default run; the same parameter pair is covered by the second integrator, which treats the trough in closed form. No failure of a mathematical statement occurred in any check.

\subsection{Missingness}\label{num:sec:missing}
The following were not done, or are not covered by the run.
\begin{itemize}
\item The revised scalar audit provides outward rational interval bounds for 55 endpoint comparisons and the strengthened $3/4+q(1000)<0.83$ test. It does not enclose trajectories, verify grid-monotonicity reductions or formalise the whole proof. There is no interval-certified orbit or all-parameter computer-assisted proof. In particular the statements $\lambda_c<0$ and ``one cycle'' at the computed parameters are floating-point observations.
\item The computations cover finitely many parameters: $\mu$ from $10^{-3}$ to $10^{3}$, $\rho\le300$ in the cartography, and narrower ranges in the other checks, as stated above. The region between the Hopf collar and the large-gain threshold is sampled, not exhausted.
\item Of the cartography, the run recomputes a subsample. The structural tests on the attractor disc and on the grid cycles (\texttt{stests\_summary.csv}, \texttt{cycle\_tests.csv}), on which several rows of Table~\ref{neg:tab:candidates} rest, are shipped as tables; the run checks statements read off them and does not recompute them. The formal limit map of Section~\ref{wed:sec:relax} is likewise shipped as a table.
\item The research scripts and outputs cited by name in Sections~\ref{sec:wedge} and \ref{sec:negative}, other than those listed in Section~\ref{num:sec:results}, are shipped under \path{provenance/research_lanes/} as they were executed. They contain absolute paths, are not run by \texttt{run\_all.py}, and were not re-executed for this section. Binary caches of the research phase, superseded runs, and the downloaded literature are not shipped; the file \texttt{SKIPPED.md} in that directory lists them.
\item The searches for a second cycle are time-boxed and their random streams depend on timing, so they are repeatable in kind and not bit for bit.
\item Two tests are restricted by conditioning: the finite-difference test of $\nabla h$ is not applicable near the Hopf curve, and the rank-one identity is not counted for the test function $y$.
\item No specialist has refereed the proofs, and priority has not been established beyond the literature search described in Section~\ref{sec:model}.
\end{itemize}
